\subsection{What did not work}
\label{sec:negative}
\textbf{Claim tested:} control of a decision is enough to solve more tasks. It is not, for every direction other than explore-more. Appendix~\ref{app:negative} gives each attempt with its numbers. Solved counts are out of 170 development tasks, with one run per arm.

Submit-sooner moved behaviour as predicted. It solved 33, 26 and 27 at three strengths, against 36 for the RL adapter it was written on. Think-briefly solved 28, 28 and 25 against 29. A test-before-submitting edit on the RL adapter solved 17 against 30 for that adapter. A direction for ``the patch being submitted is right'' reaches a held-out AUC of 0.61--0.62, against 0.77--0.83 for the explore direction by the same method. Edits aimed at failed file edits left the share of well-formed regenerated edits unchanged (Appendix~\ref{app:negative}). Two written rules gave no gain as prompt text: 25, 29 and 30 for the first rule, the second and both, against 30 for the unedited model. Compiling them into weights reproduced 53\% and 41\% of their activation shift at rank 16, with no solve-rate gain. A one-shot weight deletion chosen by gradient scores without weight updates (after \citealp{Lu2026SurprisingEffectiveness}) scored below random deletions on one batch of 64 tasks.

A pre-registered screen of 18 further edits and two random-direction controls built on the loss after long exploration (Appendix~\ref{app:round4}) found no edit that met its decision rule. Its best edits averaged $+5.0$ over their two batches, against the required $+8$. A random-direction control gained 7 in one batch, and one depth variant broke the model.

\textbf{Prediction ledger.} Decision-level predictions held for submit-sooner, think-briefly and explore-more. They failed for the failed-edit neurons. Solve-rate predictions were right in sign only for explore-more on development tasks.


\subsection{Behavioural predictions across all rounds}
\label{sec:behav}
\textbf{What the author sees.} The LLM author never reads weight values. Its inputs are behaviour statistics over 2,210 unedited episodes and activations stored at labelled decision points. It forms a class-mean-difference direction (``the next call makes the first edit'') and fits a ridge gate (held-out AUC 0.99 and 0.95 for the two gates of \texttt{pe\_lo}, the round-4 low-dose explore edit that combines the before-first-edit and early-calls gates; Appendix~\ref{app:round4}). It then writes one quiet MLP neuron per layer. It solves for the write size so that the recomputed shift on stored activations equals a target. Base weights enter only as inputs to that computation: they choose a quiet neuron and undo the output normalisation.

\textbf{Scorecard.} We scored every pre-stated behavioural prediction of the benchmark and of rounds 4 and 5 against the unedited arm of the same batch. The measure is calls before the first file edit: its mean and the share of episodes with 8 or more. The predicted direction held for 18 of 30 edits (one-sided sign test $p = 0.18$). The 30 are the edits whose prediction stated a direction for calls before the first edit and that the scorecard tabulates (\texttt{interp/REPORT.md}). Six edits are not counted: Claude's early-calls gate, whose prediction was a ceiling on the 8+ share rather than a direction (it rose from 16.0\% to 29.8\%); GPT's re-check edit, whose prediction concerned another measure that was not recorded; and four round-4 edits that the scorecard does not tabulate (\texttt{pre\_early}, \texttt{late30}, \texttt{late60}, \texttt{pu}). The record splits by edit type. Twelve edits are single-direction ``explore first'' pushes: Claude's three benchmark explore edits (\texttt{xp\_const}, always-on 0.125; \texttt{xp\_pre}, before-first-edit gate; \texttt{xp\_think}, generated-positions gate) and nine round-4 variants of the before-first-edit gate (\texttt{pe\_lo}, \texttt{pe\_hi}, \texttt{pe\_L5}, \texttt{pe\_L14}, \texttt{pe\_L20}, \texttt{pe\_L22}, \texttt{pe\_c6}, \texttt{pre15}, \texttt{pre35}; the \texttt{pe*} edits add the early-calls gate along the same direction). Five are round-4 composites (\texttt{u33}, \texttt{u36}, \texttt{u63}, \texttt{u33\_L5}, \texttt{u33\_c6}). The other 13 are GPT's three and Gemini's four scored benchmark edits and the six round-5 commit-late edits; 7 of these 13 held (4 benchmark, 3 round-5). The explore pushes moved the decision as predicted in 11 of 12 edits ($p = 0.003$). The 12 share one direction, so this sign test is descriptive. They raised mean calls before the first edit in all 46 of their edit--batch comparisons (not independent), excluding the collapsed layer-14 variant. A larger dose moved behaviour more in 4 of 5 dose pairs. Five round-4 edits combined an early explore push with a late commit push, and we predicted that they would cut episodes with 8 or more calls before the first edit. All five raised that share instead, in 25 of 25 batches. The likely reason is structural timing. The late gate fires only once 8 calls have passed without a file edit (10 in \texttt{u33\_c6}), so it cannot stop an episode from entering the 8+ bin, while the early push delays the first edit. Alone, the commit-late edits \texttt{late30} and \texttt{late60} lowered the 8+ share by only 0.013 and 0.045 in round 4. A second explanation is that shifts measured on unedited-model states do not compose once the edited rollout changes the states that the late gate sees. Of six random-direction controls, three stayed flat. Claude's benchmark control raised mean calls before the first edit by 0.45. The two round-4 controls left the median unchanged and raised the 8+ share by 0.05--0.06.

\textbf{Specificity and a warning sign.} The pushes also raise total calls per episode (for example 15.2 against 11.1 for the high-dose explore gate). The one collapse, the layer-14 variant, made the agent submit after a median of one call. In that variant, 42\% of turns had no tool call and 24\% of calls were malformed. The collapse was visible before serving. At layer 14 the write-size solver hit its ceiling, and the recomputed shift changed sign across position groups. Every working edit had small write sizes and a shift of one sign. Behavioural shift did not track solve-rate change across edits (rank correlation about zero over 16 edits). The claim is therefore direction-level behavioural control by single-direction edits, not outcome prediction. (Scripts and per-edit tables: \texttt{interp/}.)

\textbf{Tasks the unedited model never solves.} We pooled all saved development runs with at least 160 finished tasks: 42 unedited, 48 random-control, 194 edit and 5 RL-adapter runs. The unedited model never solved 108 of the 170 tasks. Edited, control or RL runs solved 28 of those at least once. Per run, solves of such tasks average 0.14 for an unedited run (counting tasks no other unedited run solved), 0.52 for a random control, 0.59 for an LLM-written edit and 2.80 for the RL adapter. Edits reach tasks that the unedited model does not solve, but random controls of the same size do so almost as often. Most of this reach therefore comes from perturbation, not from the edit's direction. The RL adapter is the exception. This analysis is post hoc, and runs of one edit are not independent (\texttt{novel\_split.py}, \texttt{novel\_passk.py}).
